\chapter{Related Work}
\label{ch:related}

This chapter reviews five groups of earlier work that are closely related to this thesis: psychiatric LLM diagnosis, retrieval with similar examples, medical multi-agent systems, criterion-level records that can be checked, and primary-diagnosis selection. The goal is not to rank the systems. The goal is to show what part of diagnosis each system studies, what outputs it keeps, and what questions are still open.

\section{Psychiatric LLM Diagnosis}

Psychiatric LLM studies cover many different tasks. Some test psychiatric knowledge or licensing questions. Others use short clinical cases, focus on one disorder, detect depression from text, diagnose from a fixed record, or support an interactive consultation~\citep{qiu2023largeai,li2024taiwan,raballo2025,sarma2025,li2025addreview,mcduff2025amie}. These tasks are not the same. They differ in the evidence given to the model, the number of possible diagnoses, the required output, and the reference used for scoring.

Screening is also different from choosing among several similar diagnoses. A screening system may only need to find a broad depression or anxiety signal. A differential-diagnosis system must compare disorders that share symptoms and still choose one primary diagnosis. This choice may depend on symptom duration, which symptoms started first, illness course, other possible causes, and comorbidity. LingxiDiagBench shows this difference: performance falls as the task moves from broad depression--anxiety classification to comorbidity recognition and twelve-category diagnosis~\citep{xu2026lingxidiagbench}. MDD-5k is another synthetic Chinese dialogue resource generated from anonymized patient cases~\citep{yin2024mdd5k}. Both datasets support repeatable experiments, but neither replaces testing on natural clinical conversations.

Another key difference is whether the system can ask more questions. MIND, WiseMind/ProAI, MAGI, DSM5AgentFlow, and conversational diagnostic AI study interactive workflows that can ask questions during the assessment~\citep{li2026mind,wu2026wisemind,bi2025magi,ozgun2025dsm5agentflow,tu2025conversational}. MentalSeek-Dx and MoodAngels instead work with information that is already present in a record or transcript~\citep{sun2026mentalseek,xiao2025moodangels}. These settings test different abilities. An interactive system can collect new evidence. A fixed-input system must work with the same possibly incomplete record. This thesis uses the fixed-input setting, so the model receives no credit for asking follow-up questions.

Most published studies still focus on the final diagnosis. A final score does not show whether a reference diagnosis was never proposed, failed a criterion rule, or remained available but was not selected as primary. This gap motivates a stage-wise evaluation instead of final-label accuracy alone.

\section{Retrieval with Similar Cases}

Retrieval can have two different roles. Knowledge retrieval adds documents, guidelines, or other factual sources~\citep{lewis2020rag}. Example retrieval selects labeled examples or similar cases and places them in the prompt for in-context learning~\citep{liu2022incontext}. This thesis uses example retrieval. The retrieved cases come from the training source. They are examples for the model, not medical evidence about the current patient.

Similar examples can show common wording, the expected output format, possible labels, and common symptom--diagnosis patterns. However, in-context learning depends strongly on which examples are selected~\citep{liu2022incontext}. A retrieved case describes another patient or synthetic case. It may guide the model, but it does not show that the current patient has the same diagnosis.

Retrieval can also hurt performance. The examples reflect the wording and label balance of the source corpus. Frequent classes may appear many times. Very similar wording may encourage shortcuts, while a weakly related example may distract the model. Retrieval should therefore be tested as one complete setting rather than assumed to improve diagnosis.

A label-aware policy may show a wider range of diagnoses and reduce repeated examples from frequent classes. However, balancing also changes the number, order, similarity, total length, and mix of examples. A comparison between global similarity retrieval and a balanced policy is therefore a comparison between two complete retrieval settings. It is not a clean test of label balance alone.

Earlier work shows that example choice matters for in-context learning, but it does not show which retrieval policy is best for fine-grained psychiatric diagnosis. Retrieval policies should therefore be compared within the same architecture and scoring contract.

\section{Medical Multi-Agent Systems}

A medical multi-agent system uses several model calls and gives them different roles. One call may propose a diagnosis, another may extract evidence, another may criticize an answer, and another may make the final decision. Other systems use consultation, debate, voting, or consensus. MedAgents uses several role-based medical experts, while MDAgents changes the collaboration plan according to estimated task difficulty~\citep{tang2024medagents,kim2024mdagents}. Multi-agent debate is a related method in which agents exchange and revise competing answers~\citep{du2024debate}.

These designs may help. A narrow role can focus on one task. A second model call may find missing evidence. A checker may compare a diagnosis with criteria, and a structured workflow can save intermediate outputs for review. These are possible benefits, not guarantees. More agents or more discussion rounds do not automatically lead to better diagnostic accuracy.

Multi-agent systems can also repeat the same mistake. Agents that use the same model family, similar prompts, and the same evidence may fail in similar ways. A wrong candidate or an unsupported evidence statement can pass into later steps. Free-form discussion may change several parts of the workflow at once, and the last component may still choose the wrong answer even when a useful alternative was mentioned earlier.

This makes the source of a performance change hard to identify. If the candidate generator, evidence record, discussion process, and final rule all change together, a higher or lower final score is a result of the whole system. It does not show which part caused the change. A comparison with a Single LLM may also change the prompt, context, number of model calls, and compute.

For this reason, candidate generation, criterion checking, and final commitment should be saved and evaluated separately. This does not make the components independent, but it shows more clearly where the recorded disagreement appears.

\section{Auditability and Criterion-Level Records}
\label{sec:auditability-definition}

An explanation that sounds clear may still fail to match the factors that led to the answer~\citep{turpin2023unfaithful}. This thesis uses four related terms. \emph{Explainability} means the account shown to a person. \emph{Interpretability} means how easy the decision process is to understand. \emph{Faithfulness} asks whether the account matches the factors that produced the output. \emph{Auditability} means that an external record is saved and can be checked or replayed. These properties are different and should be evaluated separately in healthcare~\citep{rudin2019stop,huang2024xiai}.

Diagnostic-chain and symptom-based methods can show intermediate diagnoses, confidence values, or clinically motivated features~\citep{chen2025cod,lee2026interpretable,lyu2026depllm}. Psychiatric systems may also connect outputs to guidelines, structured interviews, symptom-to-criterion links, and supporting or conflicting evidence~\citep{li2026mind,bi2025magi,ozgun2025dsm5agentflow}. Constraint-logic methods make the rule layer easier to inspect by turning diagnostic criteria into explicit programs~\citep{kim2025clp}.

These records are useful only when the criteria and evidence links are reasonable. A model may write an unsupported evidence note. A criterion may be represented poorly, and a visible rule may still contain a weak clinical assumption. Expert review of model-generated diagnostic rules is still needed~\citep{kim2025clp}. Criterion-level outputs are therefore model records that can be reviewed. They are not clinician-validated evidence or clinical truth.

An auditable system should keep the source evidence, criterion states, rule results, intermediate diagnoses, and final decision. Even then, the record does not reveal hidden model reasoning. A criterion layer may be used only for review, or it may also affect final selection. These two roles should be tested separately because good evidence recording or high compatibility inclusion does not guarantee a good primary diagnosis.

The remaining gap is therefore not only whether a system keeps intermediate records. The records must also be evaluated together on the same cases: Was the diagnosis available as a candidate? Did it enter the compatible set? Was it finally committed as primary?

\section{Primary Diagnosis Selection}

Criterion compatibility and primary-diagnosis selection answer different questions. Compatibility asks whether one diagnosis fits the available transcript under a defined rule. Because psychiatric disorders share symptoms, more than one diagnosis may pass its own rule. Compatibility alone does not show whether a diagnosis should be primary, comorbid, secondary, or only a differential alternative.

Structured medical knowledge can help organize evidence. MedKGI combines a medical knowledge graph, guided questions, and a structured dialogue state to seek useful information during diagnosis~\citep{wang2025medkgi}. DKEC builds knowledge graphs from outside medical sources and uses their links for multilabel diagnosis prediction~\citep{ge2024dkec}. These studies show how structured information can help collect evidence or predict labels. Their tested goals are different from choosing one primary diagnosis among several criterion-compatible psychiatric candidates in a fixed transcript.

A primary selector may need to compare the candidates directly. Useful information may include symptom timing, duration, illness course, past episodes, other possible causes, which syndrome explains the case best, and whether more than one disorder is present. The reviewed methods can organize evidence or find plausible labels, but they do not directly test the case in which a benchmark-reference diagnosis remains ranked and criterion-compatible while another diagnosis is committed as primary.

This gap motivates three separate measures: genuine Top-3 candidate coverage, benchmark-reference inclusion in the criterion-compatible set, and committed-primary agreement.

\section{Research Gap and Research Questions}

The reviewed work shows different parts of the diagnostic process. Some systems produce possible diagnoses. Some save criterion-related evidence. Others mainly return one final diagnosis. Table~\ref{tab:related-work-matrix} summarizes the outputs reported in representative studies. It describes the published evaluation of each study. It does not rank overall system quality or clinical ability.

\begin{table}[htbp]
\centering
\caption{Representative psychiatric diagnostic systems and the outputs available for stage-wise disagreement analysis.}
\label{tab:related-work-matrix}
\small
\begin{adjustbox}{max width=\textwidth}
\begin{tabular}{|>{\raggedright\arraybackslash}p{2.5cm}|>{\raggedright\arraybackslash}p{2.7cm}|>{\raggedright\arraybackslash}p{3.1cm}|>{\raggedright\arraybackslash}p{4.3cm}|>{\centering\arraybackslash}p{2.0cm}|}
\hline
Study or system & Evidence setting & Ranked candidates & Criterion or evidence record & Stage-wise evaluation\\
\hline
MIND~\citep{li2026mind} & Interactive consultation & No ranked evaluation reported & Criterion-based rationales and process rewards & No\\
\hline
WiseMind/\allowbreak ProAI~\citep{wu2026wisemind} & Interactive consultation & No ranked evaluation reported & Structured DSM-guided action states & Partial\\
\hline
MAGI~\citep{bi2025magi} & Interactive MINI-guided interview & No ranked evaluation reported & Symptom and dialogue evidence linked to criteria & Partial\\
\hline
DSM5Agent\allowbreak Flow~\citep{ozgun2025dsm5agentflow} & Interactive questionnaire & No ranked evaluation reported & Utterance-level supporting and conflicting criteria & Partial\\
\hline
MentalSeek-Dx~\citep{sun2026mentalseek} & Fixed clinical record & Possible diagnoses generated; ranking not evaluated & Criteria-guided staged reasoning & Partial\\
\hline
Mood\allowbreak Angels~\citep{xiao2025moodangels} & Fixed records and assessment scales & No ranked evaluation reported & DSM-5 retrieval and symptom matching & Partial\\
\hline
Constraint-logic CDSS~\citep{kim2025clp} & Structured patient facts & None & Inspectable diagnostic rules & No\\
\hline
\end{tabular}
\end{adjustbox}

\vspace{0.4em}
\begin{minipage}{0.98\textwidth}
\footnotesize
\textit{Note.} ``Partial'' means that a study reports intermediate diagnostic or criterion-related outputs but does not separately evaluate candidate detection, criterion checking, and primary commitment on the same cases. ``Not reported'' refers only to the published evaluation and does not show that the implementation lacks the capability.
\end{minipage}
\end{table}
\FloatBarrier

The reviewed work leaves five questions that are directly related to this thesis:

\begin{enumerate}
    \item Conventional classifiers, Single LLMs, and multi-stage systems often use different prompts, model calls, retrieval settings, and output contracts. A final score therefore compares complete systems, not one isolated component.
    \item Final-label accuracy does not show whether a reference diagnosis was missing from the candidates, absent from the criterion-compatible set, or available in both views but not committed as primary.
    \item Example retrieval depends on which examples are selected. Global similarity retrieval and label-aware balancing should therefore be compared within each architecture.
    \item Multi-agent and criterion-based selection methods often change several parts of the workflow at once. This makes it hard to claim that one final-selection component caused the result.
    \item A stage-wise analysis should be checked on a second synthetic source. One corpus cannot show whether the same output views remain useful after the data source changes.
\end{enumerate}

HiED addresses these questions by saving the ranked candidates, criterion states, criterion-compatible set, and final commitment under clear comparison rules.

This thesis uses the following research questions:

\begin{description}[style=nextline,leftmargin=1.5cm,labelwidth=1.1cm,itemsep=0pt,parsep=0pt,topsep=0pt,partopsep=0pt]
\item[RQ1] How does HiED perform relative to conventional classification and a Single LLM on the same internal split and parent-label evaluation?
\item[RQ2] Where is benchmark disagreement recorded across candidate generation, criterion checking, compatibility analysis, and primary commitment?
\item[RQ3] How do similar-case retrieval strategies affect Single and HiED?
\item[RQ4] Do the tested primary-selection strategies reduce the recorded selection disagreement?
\item[RQ5] Does the stage-wise analysis remain useful under a second synthetic source?
\end{description}

RQ3 first compares retrieval settings on the public-validation split and freezes the selected settings. RQ1 then tests those settings on the fixed internal held-out split. RQ2 uses the internal HiED trace to locate recorded disagreement. RQ4 compares the tested primary-selection methods. RQ5 applies the same output views to MDD-5k and uses the cross-corpus results to define the synthetic-source boundary.
